% 312_Anhang_Altermagnet_En_ch.tex
% Appendix A to Doc. 312 (Closure Scale / Lorentz section)
% Altermagnetism as a laboratory analogy: local ≠ global in symmetry breaking
% August 2026 — Johann Pascher

\providecommand{\B}{}\renewcommand{\B}{\textbf{[B]}}
\providecommand{\K}{}\renewcommand{\K}{\textbf{[K]}}
\providecommand{\E}{}\renewcommand{\E}{\textbf{[E]}}
\providecommand{\SET}{}\renewcommand{\SET}{\textbf{[AXIOM]}}

\definecolor{ffgftblue}{RGB}{30,90,160}
\definecolor{proofgreen}{RGB}{0,110,50}
\definecolor{warnviolet}{RGB}{100,0,130}
\tcbset{
  base/.style={breakable,enhanced,boxrule=0.6pt,arc=3pt,
    fonttitle=\bfseries\small,coltitle=white,top=4pt,bottom=4pt,left=6pt,right=6pt},
  ffgftbox/.style={base,colback=blue!5,colframe=ffgftblue,colbacktitle=ffgftblue},
  proofbox/.style={base,colback=green!5,colframe=proofgreen,colbacktitle=proofgreen},
  warnbox/.style={base,colback=violet!6,colframe=violet!60!black,colbacktitle=violet!60!black},
}

% ============================================================
\chapter*{Appendix A to Doc.\,312 — Altermagnetism as a Laboratory Analogy}
\addcontentsline{toc}{chapter}{Appendix A: Altermagnetism as a Laboratory Analogy}
% ============================================================

\section*{Purpose of This Appendix}

Doc.\,312, §\,Lorentz (source lines 936--947), makes the statement:

\begin{quote}
\itshape
On compact topology Lorentz invariance holds locally without restriction,
but boosts are no longer global symmetries.
The distinction is topological; no local measurement can detect it —
only global markers.
\end{quote}

The concept \enquote{locally broken, globally preserved} is abstract and
difficult to ground without experimental reference. Altermagnetism
provides a directly measured, mathematically fully analysed example
of the same structure at a completely different scale. This appendix
documents the analogy and marks its limits.

\textbf{Purpose:} Didactic reference and terminological calibration
for \enquote{local $\neq$ global}. No new physical content for FFGFT.
Source: Šmejkal, Sinova, Jungwirth,
\textit{Emerging Research Landscape of Altermagnetism},
Phys.\ Rev.\ X \textbf{12}, 040501 (2022).

% ============================================================
\section{Altermagnetism: Three-Way Classification via Spin Groups \E}
% ============================================================

The non-relativistic spin-group classification of collinear magnets
yields exactly three phases, determined by which symmetry operation
connects the magnetic sublattices:

\begin{center}
\begin{tabular}{lllc}
\toprule
\textbf{Type} & \textbf{Sublattice link} & \textbf{Kramers deg.} & $M_{\rm net}$\\
\midrule
Ferromagnet     & none (one spin sublattice)         & lifted (Zeeman)   & $\neq 0$\\
Antiferromagnet & translation $t$ or inversion $P$   & preserved         & $= 0$\\
Altermagnet     & rotation $C_n$ ($n=2,4,6$)         & locally lifted    & $= 0$\\
\bottomrule
\end{tabular}
\end{center}

The decisive distinction between antiferromagnet and altermagnet: in
both, the net magnetisation is zero. But in the antiferromagnet a
translation or inversion connects the sublattices — this
\emph{enforces} Kramers degeneracy
$E(\mathbf{k},\uparrow) = E(\mathbf{k},\downarrow)$ across the entire
Brillouin zone. In the altermagnet a genuine rotation ($C_4$ in RuO$_2$)
connects the sublattices — this lifts the degeneracy in a
direction-dependent manner.

% ============================================================
\section{The Model Hamiltonian and Its Symmetries \E}
% ============================================================

The effective single-particle Hamiltonian of the $d_{xy}$-wave
altermagnet (Šmejkal et al., Eq.\ 3):
\begin{equation}
  H(\mathbf{k}) = 2t\cos k_x \cos k_y \cdot \sigma_0
    + 2t_J \sin k_x \sin k_y \cdot \sigma_z.
  \label{eq:alt_ham}
\end{equation}
The energy eigenvalues are
$E_\pm(\mathbf{k}) = 2t\cos k_x \cos k_y \pm 2t_J \sin k_x \sin k_y$
with spin splitting
\begin{equation}
  \Delta E(\mathbf{k}) = E_+ - E_- = 4t_J \sin k_x \sin k_y.
  \label{eq:splitting}
\end{equation}

\begin{tcolorbox}[proofbox,title={Proof: $C_4$ maps the spin-up onto the spin-down Fermi surface \B}]
The $C_4$ rotation in $\mathbf{k}$-space acts as
$(k_x, k_y) \to (-k_y, k_x)$. Then:
\begin{align}
  E_{\rm dn}(-k_y,\, k_x)
  &= 2t\cos(-k_y)\cos(k_x) - 2t_J\sin(-k_y)\sin(k_x) \notag\\
  &= 2t\cos(k_y)\cos(k_x) + 2t_J\sin(k_y)\sin(k_x)
   = E_{\rm up}(k_x, k_y). \quad\checkmark
\end{align}
The spin-up Fermi surface is mapped exactly onto the spin-down surface
by $C_4$ (numerically confirmed, \texttt{pruef\_altermagnet\_1\_symmetrie.py}).
\end{tcolorbox}

Two further analytically provable statements:

\begin{enumerate}
\item \textbf{Global compensation:}
  $\displaystyle\frac{1}{(2\pi)^2}
  \int_{\mathrm{BZ}} \Delta E(\mathbf{k})\,\mathrm{d}^2k = 0$,
  since $\int_{-\pi}^{\pi}\sin k_x\,\mathrm{d}k_x = 0$.
  No macroscopic net magnetisation. \B

\item \textbf{Nodal lines:}
  $\Delta E = 0$ exactly at $k_x = 0$ or $k_y = 0$ (two orthogonal
  nodal lines, protected by mirror-plane symmetries that map one spin
  sublattice onto the other). \B
\end{enumerate}

% ============================================================
\section{Locally Broken, Globally Preserved — the Shared Structure}
% ============================================================

\begin{tcolorbox}[ffgftbox,title={Structural comparison with Doc.\,312 §\,Lorentz}]
\begin{tabular}{p{0.45\linewidth}p{0.45\linewidth}}
\textbf{FFGFT / Doc.\,312} & \textbf{Altermagnet / Šmejkal 2022}\\[4pt]
Compact topology $T^4/\mathbb{Z}_3$ & Crystal lattice with $C_4$ rotation\\[2pt]
Lorentz invariance holds locally & $T$-symmetry locally broken ($\Delta E \neq 0$)\\[2pt]
Boosts: no longer global symmetries & Spin: globally compensated ($M=0$)\\[2pt]
Only topological markers detectable & Only ARPES/spin-current measurable\\[2pt]
Scale: Planck length & Scale: ångström (crystal lattice)
\end{tabular}
\end{tcolorbox}

The structure is the same: a symmetry (time-reversal $T$ or Lorentz
boost) is locally broken in the sense that it fails at a single point
(in $\mathbf{k}$-space or for a specific direction of motion in space).
But the global integral — over the entire Brillouin zone or over all
directions — is zero: the overall structure is compensated. Locally
this breaking is detectable only by direction-resolved measurements;
an isotropic probe reveals nothing.

% ============================================================
\section{Limits of the Analogy \E}
% ============================================================

\begin{tcolorbox}[warnbox,title={Restrictions — explicit}]
\begin{enumerate}
  \item \textbf{Different scales.} Altermagnetism: $\sim$\,\AA ngström.
    FFGFT: Planck scale ($\sim 10^{-35}$\,m). No shared observables.
  \item \textbf{Different symmetry groups.} The spin group $r \times R_s$
    of the altermagnet and the orbifold group $T^4/\mathbb{Z}_3$
    of FFGFT are different mathematical objects.
  \item \textbf{No information flow.} FFGFT does not explain solid-state
    effects. Altermagnetism contributes no new derivations to FFGFT.
  \item \textbf{Direction of the analogy.} The analogy serves the
    didactic calibration of the concept \enquote{locally $\neq$ globally}.
    It is not a derivation and must not be read as a confirmation of FFGFT.
\end{enumerate}
\end{tcolorbox}

% ============================================================
\section{Verification}
% ============================================================

All statements in §\,2 are numerically and symbolically confirmed by
\texttt{pruef\_altermagnet\_1\_symmetrie.py}. The script is located in
\texttt{2/python/Dok312\_Anhang\_Skripte/}.

Reference: L.\ Šmejkal, J.\ Sinova, T.\ Jungwirth,
\textit{Emerging Research Landscape of Altermagnetism},
Phys.\ Rev.\ X \textbf{12}, 040501 (2022).
arXiv:2204.10844.
